%% kapa: exam=1 round=28 year=2017 period=1 form=A subject=economics subject_ko=경제학원론 seq=31 number=71
%% source: Q-Net 감정평가사 기출문제 (cst00309) artlSeq=5199603 file=extracted/1008/5199603/2017년 제28회 감정평가사 1차시험 문제지/1교시 A형.pdf
%% conversion: kordoc 4.24.51 -> tools/gen1.py (auto); review=proofread-sonnet; flags=none
\begin{question}{71}
다음 폐쇄경제 IS-LM 모형에서 경제는 균형을 이루고 있고, 현재 명목화폐공급량($M$)은 2이다. 중앙은행은 확장적 통화정책을 실시하여 현재보다 균형 이자율을 0.5만큼 낮추고, 균형국민소득을 증가시키고자 한다. 이를 위한 명목화폐공급량의 증가분($\Delta M$)은? (단, $Y$는 국민소득, $r$은 이자율, $M^{d}$는 명목화폐수요량, $P$는 물가이고 1로 불변이다.)
\begin{framedtext}
○ IS 곡선: $r=4-0.05Y$ \par
○ 실질화폐수요함수: $\frac{M^{d}}{P}=0.15Y-r$
\end{framedtext}
\begin{choices}
  \item 0.5
  \item 2
  \item 2.5
  \item 3
  \item 4
\end{choices}
\end{question}
\answer{2}
